\documentclass[11pt]{article}

\usepackage[margin=25mm]{geometry}

\usepackage{amsmath,amssymb,fontspec,unicode-math}

\IfFontExistsTF{DejaVu Serif}{\setmainfont{DejaVu Serif}}{\setmainfont{Latin Modern Roman}}

\IfFontExistsTF{Latin Modern Math}{\setmathfont{Latin Modern Math}}{\setmathfont{DejaVu Math TeX Gyre}}

\usepackage{xurl,hyperref}

\setlength{\parindent}{0pt}\setlength{\parskip}{6pt}

\begin{document}

\title{Beyond Õ(√T) Constraint Violation for Online Convex Optimization with Adversarial Constraints}

\author{ELVIA preliminary AI analysis}\date{21 October 2026}\maketitle

Preliminary assessment: Major.\par

Independent scientific review remains pending. Author adoption has not been requested.

The smooth small-loss bound lacks a domain/self-bounding hypothesis: an exact one-round convex loss satisfies local nonnegativity and smoothness but violates the displayed regret constant. The expert and covering arguments are separate and remain unrefuted.

\section*{Review scope}

Claim-focused reading of the original paper and relevant author-released passages; direct mathematical or static-source reasoning. No author program, earlier verifier, model or benchmark was executed during this check.

\section*{Problem 1: Local nonnegativity does not imply the small-loss bound}

Source: Theorem 4, Eq.(22), PDF p.9; Assumptions 1,4, PDF pp.2,9; Eq.(5), PDF p.4

Assessment: supported; impact: major

\subsection*{Source claim (paraphrase)}

For nonnegative convex \(M\)-smooth losses on a closed decision set of diameter \(D\), Eq.(22) bounds regret by \(4D\sqrt{M\sum_t l_t(u)}+4D^2M\).

\subsection*{Analysis}

Take \(\mathcal X=[0,1]\), \(D=T=x_1=1\), comparator \(u=0\), and \(l(x)=x/2+x^2/200\). On the stated domain, \(0\le l\le0.505<1\), while \(l^{\prime\prime}=M=0.01\) exactly. The comparator loss is zero, so the claimed right side is \(0.04\), but the first-round regret is \(0.505\). The nonzero gradient \(0.51\) avoids step-size singularity; no update can undo an already incurred first loss. Setting the ordinary cost to zero and constraint penalty to this loss gives the same violation of the explicit constraint-only illustration. The polynomial is globally smooth but is not globally nonnegative.

\subsection*{Finding and scope}

The local-domain theorem and explicit constant fail. This does not directly refute the expert/covering route or establish a contrary asymptotic rate for every smooth extension.

\subsection*{Proposed correction}

Add \(\|\nabla l_t(x)\|^2\le2Ml_t(x)\), for example through a globally nonnegative smooth extension. Adaptive projection then gives \(R\le2D\sqrt{M(R+L)}\), whence \(R\le4D^2M+2D\sqrt{ML}\), restoring the more conservative printed constants. Weighted Cauchy–Schwarz supplies the corresponding self-bound for nonnegative surrogates.

This domain omission is inherited from the cited older monograph wording; its later edition explicitly assumes self-boundedness. No paper-specific author adoption is established.

\section*{Limits}

Preliminary AI assessment; independent human scientific adjudication remains pending.\par

Analytic witnesses and static release checks do not remeasure the published benchmark results.\par

\section*{Source identity}

NeurIPS 2025 official proceedings PDF (retained local copy)

\url{https://proceedings.neurips.cc/paper_files/paper/2025/file/d1e2165d5238dda55582fc55d193d170-Paper-Conference.pdf}

Reviewed PDF SHA-256: \nolinkurl{13ad20a8b470de0ddc328484e2f4623c912a0c31622b23d24cc693858a4ce628}

The SHA-256 identifies the exact retained PDF inspected in this rapid source/argument check. Original supplied report bytes are retained locally. No unambiguous original PDF retrieval timestamp was recovered; file modification time was not substituted.

\end{document}
